\documentclass[11pt,a4paper]{article}
\usepackage[T1]{fontenc}
\usepackage[utf8]{inputenc}
\usepackage[brazil]{babel}
\usepackage{amsmath,amssymb,textcomp}
\usepackage[margin=2.2cm]{geometry}
\usepackage{enumitem}
\usepackage{array}
\setlength{\parindent}{0pt}
\setlength{\parskip}{0.6em}

\begin{document}
\section*{Avaliação de questões --- avaliador 1}
\subsection*{Como responder}

São 22 questões de Matemática do Ensino Médio. Leia cada uma como leria uma
questão que você pensasse em usar com a sua turma, e responda às quatro
perguntas que vêm logo abaixo dela.

Onde se pergunta pelo \textbf{nível de dificuldade}, responda pensando nos
seus alunos, não num aluno ideal. É essa resposta que permite calibrar o
sistema: ele classifica cada questão como fácil, média ou difícil por conta
própria, e precisamos saber o quanto isso corresponde à sala de aula real.

\textbf{A amostra tem qualidade variável, e pode conter questões com erro
matemático.} Se encontrar algum, aponte onde.

Não há resposta certa sobre você: o que está sendo avaliado é o sistema que
produziu estas questões, não quem as lê. Um ``recusada'' bem justificado vale
mais para este trabalho do que um ``aceita'' por gentileza.

Você pode responder direto neste documento impresso ou na planilha
\texttt{respostas-avaliador-1.csv}, que já vem com os códigos dos itens nas linhas.

Tempo estimado: cerca de 75 minutos.
\bigskip\hrule\bigskip
\begin{samepage}
\subsection*{Q012}
Uma companhia de saneamento cobra a conta mensal de água residencial segundo a tabela abaixo, válida apenas para consumos entre 0 e 35 m³ por mês, pois essa é a capacidade máxima de registro do hidrômetro residencial utilizado:

\begin{itemize}[nosep,leftmargin=*]
\item Consumo de 0 até 10 m³: tarifa mínima fixa de R\$\,25,00, independente do volume efetivamente consumido.
\item Consumo acima de 10 m³ e até 25 m³: R\$\,25,00 mais R\$\,3,50 por m³ que exceder 10 m³.
\item Consumo acima de 25 m³ e até 35 m³: o valor cobrado na faixa anterior no limite de 25 m³ (ou seja, R\$\,77,50) mais R\$\,5,00 por m³ que exceder 25 m³.
\end{itemize}

Seja $x$ o consumo mensal em metros cúbicos, com $0 \le x \le 35$, e $C(x)$ o valor da conta, em reais, correspondente a esse consumo.

Qual das alternativas a seguir descreve corretamente o gráfico de $C$ e a imagem dessa função no domínio de validade dado?
\begin{enumerate}[label=(\alph*),nosep,leftmargin=*]
\item O gráfico é um segmento horizontal em $C=25$ para $0 \le x \le 10$; em seguida uma semirreta de coeficiente angular $3{,}5$ unindo $(10;25)$ a $(25;77{,}5)$; e por fim uma semirreta de coeficiente angular $5$ unindo $(25;77{,}5)$ a $(35;127{,}5)$. A função é constante em $[0,10]$ e crescente em $(10,35]$, com imagem $[25;\,127{,}5]$.
\item A função é crescente em todo o domínio $[0,35]$, sem nenhum trecho constante, já que o valor da conta aumenta proporcionalmente ao consumo desde $x=0$; a imagem é $[25;\,127{,}5]$.
\item O gráfico é constante em $[0,10]$ e depois crescente, mas como a tarifa de R\$\,3,50 e a de R\$\,5,00 incidem sobre todo o consumo da faixa (e não apenas sobre o excedente), o valor máximo é $C(35) = 162{,}5$, de modo que a imagem é $[25;\,162{,}5]$.
\item Como a última faixa de cobrança (acima de 25 m³) não tem um valor de tarifa diferente após um certo ponto, o consumo pode crescer indefinidamente e a imagem da função é o intervalo ilimitado $[25, +\infty)$.
\end{enumerate}

\textbf{Gabarito proposto:} Alternativa A — $C$ é constante ($=25$) em $[0,10]$, crescente em $(10,35]$ (inclinação $3{,}5$ até $x=25$ e depois $5$), com imagem $[25;\,127{,}5]$.

\textbf{Habilidade declarada:} EM13MAT405 --- Reconhecer funções definidas por uma ou mais sentenças (como a tabela do Imposto de Renda, contas de luz, água, gás etc.), em suas representações algébrica e gráfica, convertendo essas representações de uma para outra e identificando domínios de validade, imagem, crescimento e decrescimento.
\end{samepage}

\noindent\rule{\linewidth}{0.4pt}
\textbf{Tem erro matemático?}\quad
\mbox{$\square$~não} \quad \mbox{$\square$~sim, aqui:} \underline{\hspace{5cm}}
\quad \mbox{$\square$~não sei dizer}

\textbf{Que nível de dificuldade esta questão tem para os seus alunos?}\quad
\mbox{$\square$~fácil} \quad \mbox{$\square$~média} \quad
\mbox{$\square$~difícil} \quad \mbox{$\square$~fora do alcance da turma}

\textbf{Você usaria esta questão?}\quad
\mbox{$\square$~aceita} \quad \mbox{$\square$~aceita com ajuste} \quad
\mbox{$\square$~recusada}

\textbf{Por quê?} \emph{(obrigatório quando não for ``aceita'')}

\underline{\hspace{\linewidth}}

\underline{\hspace{\linewidth}}
\vspace{0.6em}

\begin{samepage}
\subsection*{Q003}
Uma empresa de transporte por aplicativo cobra uma tarifa composta por um valor fixo (bandeirada) somado a um valor proporcional à distância percorrida, medida em quilômetros. Um passageiro que percorreu 4 km pagou R\$\,18,00 pela corrida, e outro passageiro que percorreu 10 km pagou R\$\,36,00.

a) Determine a lei que expressa o custo $C(d)$ de uma corrida (em reais) em função da distância $d$ percorrida (em quilômetros), sabendo que essa relação é modelada por uma função polinomial de 1º grau.

b) Utilizando a lei obtida, calcule quanto pagará um passageiro que percorrer 15 km.

c) Determine a distância percorrida por um passageiro que pagou R\$\,81,00 pela corrida.

\textbf{Gabarito proposto:} C(d) = 3d + 6; custo para 15 km = R\$\,51,00; distância para R\$\,81,00 = 25 km

\textbf{Habilidade declarada:} EM13MAT302 --- Resolver e elaborar problemas cujos modelos são as funções polinomiais de 1º e 2º graus, em contextos diversos, incluindo ou não tecnologias digitais.
\end{samepage}

\noindent\rule{\linewidth}{0.4pt}
\textbf{Tem erro matemático?}\quad
\mbox{$\square$~não} \quad \mbox{$\square$~sim, aqui:} \underline{\hspace{5cm}}
\quad \mbox{$\square$~não sei dizer}

\textbf{Que nível de dificuldade esta questão tem para os seus alunos?}\quad
\mbox{$\square$~fácil} \quad \mbox{$\square$~média} \quad
\mbox{$\square$~difícil} \quad \mbox{$\square$~fora do alcance da turma}

\textbf{Você usaria esta questão?}\quad
\mbox{$\square$~aceita} \quad \mbox{$\square$~aceita com ajuste} \quad
\mbox{$\square$~recusada}

\textbf{Por quê?} \emph{(obrigatório quando não for ``aceita'')}

\underline{\hspace{\linewidth}}

\underline{\hspace{\linewidth}}
\vspace{0.6em}

\begin{samepage}
\subsection*{Q020}
A tabela a seguir apresenta pares de valores (x, y) gerados por uma mesma regra de formação:

\begin{center}
\begin{tabular}{|p{0.143\linewidth}|p{0.143\linewidth}|p{0.143\linewidth}|p{0.143\linewidth}|p{0.143\linewidth}|p{0.143\linewidth}|}
\hline
\textbf{x} & \textbf{-2} & \textbf{-1} & \textbf{0} & \textbf{1} & \textbf{2} \\
\hline
y & 1 & 3 & 5 & 7 & 9 \\
\hline
\end{tabular}
\end{center}

a) Investigue os pares (x, y) da tabela, identifique um padrão entre os valores de x e y e escreva, com suas próprias palavras, a expressão algébrica y = f(x) que gera esses valores. Justifique por escrito como você chegou a essa expressão.

b) Com base no padrão observado, classifique o tipo de função obtida (por exemplo: constante, polinomial do 1º grau, polinomial do 2º grau etc.) e justifique sua classificação.

c) Assinale a alternativa que apresenta corretamente a lei de formação encontrada.

(Os itens a e b devem ser respondidos de forma discursiva e serão corrigidos separadamente, avaliando o processo de investigação e a justificativa apresentada; o item c será corrigido de forma objetiva.)
\begin{enumerate}[label=(\alph*),nosep,leftmargin=*]
\item y = 2x + 5
\item y = 2x + 1
\item y = 2x
\item y = x² + 5
\end{enumerate}

\textbf{Gabarito proposto:} y = 2x + 5 (alternativa A)

\textbf{Habilidade declarada:} EM13MAT501 --- Investigar relações entre números expressos em tabelas para representá-los no plano cartesiano, identificando padrões e criando conjecturas para generalizar e expressar algebricamente essa generalização, reconhecendo quando essa representação é de função polinomial de 1º grau.
\end{samepage}

\noindent\rule{\linewidth}{0.4pt}
\textbf{Tem erro matemático?}\quad
\mbox{$\square$~não} \quad \mbox{$\square$~sim, aqui:} \underline{\hspace{5cm}}
\quad \mbox{$\square$~não sei dizer}

\textbf{Que nível de dificuldade esta questão tem para os seus alunos?}\quad
\mbox{$\square$~fácil} \quad \mbox{$\square$~média} \quad
\mbox{$\square$~difícil} \quad \mbox{$\square$~fora do alcance da turma}

\textbf{Você usaria esta questão?}\quad
\mbox{$\square$~aceita} \quad \mbox{$\square$~aceita com ajuste} \quad
\mbox{$\square$~recusada}

\textbf{Por quê?} \emph{(obrigatório quando não for ``aceita'')}

\underline{\hspace{\linewidth}}

\underline{\hspace{\linewidth}}
\vspace{0.6em}

\begin{samepage}
\subsection*{Q045}
Considere a progressão aritmética $(a_n)$ com primeiro termo $a_1 = 7$ e razão $r = 4$, definida para $n = 1, 2, 3, \dots$

a) Determine a lei de uma função afim $f(n)$ tal que $f(n) = a_n$ para todo $n$ natural positivo, e explicite qual é o domínio dessa função.

b) Utilizando a função $f$ obtida no item (a), calcule o vigésimo termo da progressão, isto é, $a_{20}$.

\textbf{Gabarito proposto:} f(n) = 4n + 3, com domínio $\mathbb{N}^* = \{1,2,3,\dots\}$; $a_{20} = 83$

\textbf{Habilidade declarada:} EM13MAT507 --- Identificar e associar sequências numéricas (PA) a funções afins de domínios discretos para análise de propriedades, incluindo dedução de algumas fórmulas e resolução de problemas.
\end{samepage}

\noindent\rule{\linewidth}{0.4pt}
\textbf{Tem erro matemático?}\quad
\mbox{$\square$~não} \quad \mbox{$\square$~sim, aqui:} \underline{\hspace{5cm}}
\quad \mbox{$\square$~não sei dizer}

\textbf{Que nível de dificuldade esta questão tem para os seus alunos?}\quad
\mbox{$\square$~fácil} \quad \mbox{$\square$~média} \quad
\mbox{$\square$~difícil} \quad \mbox{$\square$~fora do alcance da turma}

\textbf{Você usaria esta questão?}\quad
\mbox{$\square$~aceita} \quad \mbox{$\square$~aceita com ajuste} \quad
\mbox{$\square$~recusada}

\textbf{Por quê?} \emph{(obrigatório quando não for ``aceita'')}

\underline{\hspace{\linewidth}}

\underline{\hspace{\linewidth}}
\vspace{0.6em}

\begin{samepage}
\subsection*{Q061}
Um biólogo inicia uma cultura de bactérias com 200 indivíduos. A cada hora que passa, a quantidade de bactérias na cultura dobra em relação à hora anterior. Seja $a_n$ o número de bactérias existentes ao final da $n$-ésima hora de observação, com $n = 1, 2, 3, \dots$ (isto é, $a_1$ é a quantidade ao final da 1ª hora, $a_2$ ao final da 2ª hora, e assim por diante).

a) Mostre que a sequência $(a_n)$ é uma progressão geométrica, indicando o primeiro termo e a razão.

b) Escreva a lei $f(n)$ de uma função exponencial que descreva exatamente essa progressão geométrica, especificando qual é o domínio dessa função e por que ele deve ser desse tipo (e não todos os números reais).

c) Usando a fórmula do termo geral da PG (ou, equivalentemente, a função $f(n)$ obtida no item b), determine quantas bactérias haverá ao final da 6ª hora de observação.

\textbf{Gabarito proposto:} a) PG de razão $q=2$ e $a_1=200$; b) $f(n)=200\cdot 2^{n-1}$, com domínio $D=\{1,2,3,\dots\}$ (naturais, pois $n$ conta horas); c) $a_6 = f(6) = 6400$ bactérias.

\textbf{Habilidade declarada:} EM13MAT508 --- Identificar e associar sequências numéricas (PG) a funções exponenciais de domínios discretos para análise de propriedades, incluindo dedução de algumas fórmulas e resolução de problemas.
\end{samepage}

\noindent\rule{\linewidth}{0.4pt}
\textbf{Tem erro matemático?}\quad
\mbox{$\square$~não} \quad \mbox{$\square$~sim, aqui:} \underline{\hspace{5cm}}
\quad \mbox{$\square$~não sei dizer}

\textbf{Que nível de dificuldade esta questão tem para os seus alunos?}\quad
\mbox{$\square$~fácil} \quad \mbox{$\square$~média} \quad
\mbox{$\square$~difícil} \quad \mbox{$\square$~fora do alcance da turma}

\textbf{Você usaria esta questão?}\quad
\mbox{$\square$~aceita} \quad \mbox{$\square$~aceita com ajuste} \quad
\mbox{$\square$~recusada}

\textbf{Por quê?} \emph{(obrigatório quando não for ``aceita'')}

\underline{\hspace{\linewidth}}

\underline{\hspace{\linewidth}}
\vspace{0.6em}

\begin{samepage}
\subsection*{Q053}
Uma academia cobra R\$\,70,00 pela mensalidade do primeiro mês de plano. A partir do segundo mês, o valor da mensalidade aumenta R\$\,8,00 em relação ao mês imediatamente anterior, e esse padrão se mantém indefinidamente. Chame de $n$ o número do mês (com $n = 1, 2, 3, \dots$) e de $f(n)$ o valor da mensalidade cobrada nesse mês. Construa a lei que expressa $f(n)$ em função de $n$, levando em conta o conjunto de valores que $n$ pode assumir, e use essa lei para determinar o valor da mensalidade no 15º mês.
\begin{enumerate}[label=(\alph*),nosep,leftmargin=*]
\item R\$\,182,00
\item R\$\,190,00
\item R\$\,1.050,00
\item R\$\,120,00
\end{enumerate}

\textbf{Gabarito proposto:} R\$\,182,00

\textbf{Habilidade declarada:} EM13MAT507 --- Identificar e associar sequências numéricas (PA) a funções afins de domínios discretos para análise de propriedades, incluindo dedução de algumas fórmulas e resolução de problemas.
\end{samepage}

\noindent\rule{\linewidth}{0.4pt}
\textbf{Tem erro matemático?}\quad
\mbox{$\square$~não} \quad \mbox{$\square$~sim, aqui:} \underline{\hspace{5cm}}
\quad \mbox{$\square$~não sei dizer}

\textbf{Que nível de dificuldade esta questão tem para os seus alunos?}\quad
\mbox{$\square$~fácil} \quad \mbox{$\square$~média} \quad
\mbox{$\square$~difícil} \quad \mbox{$\square$~fora do alcance da turma}

\textbf{Você usaria esta questão?}\quad
\mbox{$\square$~aceita} \quad \mbox{$\square$~aceita com ajuste} \quad
\mbox{$\square$~recusada}

\textbf{Por quê?} \emph{(obrigatório quando não for ``aceita'')}

\underline{\hspace{\linewidth}}

\underline{\hspace{\linewidth}}
\vspace{0.6em}

\begin{samepage}
\subsection*{Q037}
Duas gráficas de impressão, identificadas por R e S, cobram pelo serviço de acordo com o número de folhas impressas ($x$) e o valor total cobrado, em reais ($y$). Usando um aplicativo de geometria dinâmica, uma estudante registrou os seguintes pontos no plano cartesiano:

\begin{itemize}[nosep,leftmargin=*]
\item O gráfico de R passa pelos pontos $A(0,3)$ e $B(4,11)$.
\item O gráfico de S passa pelos pontos $C(2,4)$ e $D(5,10)$.
\end{itemize}

Determine as leis $y = f(x)$ que descrevem R e S e, a partir delas, analise o comportamento geométrico das duas retas no plano cartesiano, comparando-as com a reta de equação $y = 2x$. Assinale a alternativa que classifica corretamente cada gráfico quanto ao tipo de função e apresenta a justificativa geométrica adequada.
\begin{enumerate}[label=(\alph*),nosep,leftmargin=*]
\item R é uma função afim não proporcional ($y=2x+3$), cujo gráfico é uma reta paralela a $y=2x$, deslocada 3 unidades para cima, não passando pela origem; S é uma função linear (proporcional), pois sua lei é $y=2x$, coincidindo com a própria reta $y=2x$ e passando pela origem do plano cartesiano.
\item Como R e S têm o mesmo coeficiente angular ($a=2$) da reta $y=2x$, ambas representam relações de proporcionalidade direta, diferindo apenas na posição inicial da reta no plano.
\item R é proporcional, pois seu gráfico foi definido a partir de um ponto com abscissa nula ($x=0$); S não é proporcional, pois nenhum dos pontos fornecidos tem abscissa igual a zero, logo sua reta não passa pela origem.
\item Nem R nem S representam relações de proporcionalidade direta, pois ambas as leis obtidas apresentam termo constante não nulo: $y=2x+3$ para R e $y=2x+2$ para S; portanto, nenhuma das retas passa pela origem.
\end{enumerate}

\textbf{Gabarito proposto:} R é uma função afim não proporcional ($y=2x+3$), cujo gráfico é paralelo a $y=2x$ mas deslocado 3 unidades para cima, não passando pela origem; S é uma função linear (proporcional), pois $y=2x$ coincide com a própria reta $y=2x$, passando pela origem.

\textbf{Habilidade declarada:} EM13MAT401 --- Converter representações algébricas de funções polinomiais de 1º grau para representações geométricas no plano cartesiano, distinguindo os casos nos quais o comportamento é proporcional, recorrendo ou não a softwares ou aplicativos de álgebra e geometria dinâmica.
\end{samepage}

\noindent\rule{\linewidth}{0.4pt}
\textbf{Tem erro matemático?}\quad
\mbox{$\square$~não} \quad \mbox{$\square$~sim, aqui:} \underline{\hspace{5cm}}
\quad \mbox{$\square$~não sei dizer}

\textbf{Que nível de dificuldade esta questão tem para os seus alunos?}\quad
\mbox{$\square$~fácil} \quad \mbox{$\square$~média} \quad
\mbox{$\square$~difícil} \quad \mbox{$\square$~fora do alcance da turma}

\textbf{Você usaria esta questão?}\quad
\mbox{$\square$~aceita} \quad \mbox{$\square$~aceita com ajuste} \quad
\mbox{$\square$~recusada}

\textbf{Por quê?} \emph{(obrigatório quando não for ``aceita'')}

\underline{\hspace{\linewidth}}

\underline{\hspace{\linewidth}}
\vspace{0.6em}

\begin{samepage}
\subsection*{Q007}
Uma oficina mecânica cobra por seus serviços de revisão um valor fixo de mão de obra, mais uma taxa que depende do número de horas de trabalho. Uma revisão que durou 2 horas custou R\$\,150,00, e outra revisão, no mesmo padrão de cobrança, durou 5 horas e custou R\$\,300,00. Sabendo que o custo total varia linearmente com o número de horas trabalhadas, qual seria o custo de uma revisão que durasse 8 horas?
\begin{enumerate}[label=(\alph*),nosep,leftmargin=*]
\item R\$\,450,00
\item R\$\,400,00
\item R\$\,600,00
\item R\$\,480,00
\end{enumerate}

\textbf{Gabarito proposto:} R\$\,450,00

\textbf{Habilidade declarada:} EM13MAT302 --- Resolver e elaborar problemas cujos modelos são as funções polinomiais de 1º e 2º graus, em contextos diversos, incluindo ou não tecnologias digitais.
\end{samepage}

\noindent\rule{\linewidth}{0.4pt}
\textbf{Tem erro matemático?}\quad
\mbox{$\square$~não} \quad \mbox{$\square$~sim, aqui:} \underline{\hspace{5cm}}
\quad \mbox{$\square$~não sei dizer}

\textbf{Que nível de dificuldade esta questão tem para os seus alunos?}\quad
\mbox{$\square$~fácil} \quad \mbox{$\square$~média} \quad
\mbox{$\square$~difícil} \quad \mbox{$\square$~fora do alcance da turma}

\textbf{Você usaria esta questão?}\quad
\mbox{$\square$~aceita} \quad \mbox{$\square$~aceita com ajuste} \quad
\mbox{$\square$~recusada}

\textbf{Por quê?} \emph{(obrigatório quando não for ``aceita'')}

\underline{\hspace{\linewidth}}

\underline{\hspace{\linewidth}}
\vspace{0.6em}

\begin{samepage}
\subsection*{Q071}
Uma confecção fabrica um certo modelo de camisa a um custo de produção de R\$\,40,00 por unidade. Um estudo de mercado mostrou que, se o preço de venda for $p$ reais por unidade, a quantidade vendida mensalmente será $q(p) = 300 - 5p$ unidades, válida para $0 < p < 60$. O lucro mensal da confecção, em reais, é dado pelo produto entre o lucro obtido em cada unidade vendida (preço de venda menos custo de produção) e a quantidade vendida naquele mês. Qual deve ser o preço de venda $p$, em reais, para que o lucro mensal da confecção seja máximo?
\begin{enumerate}[label=(\alph*),nosep,leftmargin=*]
\item R\$\,50,00
\item R\$\,30,00
\item R\$\,100,00
\item R\$\,60,00
\end{enumerate}

\textbf{Gabarito proposto:} R\$\,50,00 (alternativa a)

\textbf{Habilidade declarada:} EM13MAT503 --- Investigar pontos de máximo ou de mínimo de funções quadráticas em contextos da Matemática Financeira ou da Cinemática, entre outros.
\end{samepage}

\noindent\rule{\linewidth}{0.4pt}
\textbf{Tem erro matemático?}\quad
\mbox{$\square$~não} \quad \mbox{$\square$~sim, aqui:} \underline{\hspace{5cm}}
\quad \mbox{$\square$~não sei dizer}

\textbf{Que nível de dificuldade esta questão tem para os seus alunos?}\quad
\mbox{$\square$~fácil} \quad \mbox{$\square$~média} \quad
\mbox{$\square$~difícil} \quad \mbox{$\square$~fora do alcance da turma}

\textbf{Você usaria esta questão?}\quad
\mbox{$\square$~aceita} \quad \mbox{$\square$~aceita com ajuste} \quad
\mbox{$\square$~recusada}

\textbf{Por quê?} \emph{(obrigatório quando não for ``aceita'')}

\underline{\hspace{\linewidth}}

\underline{\hspace{\linewidth}}
\vspace{0.6em}

\begin{samepage}
\subsection*{Q052}
Uma boia de monitoramento registrou a altura da água (em metros) em um cais, em intervalos regulares de 3 horas, ao longo de um dia. Observou-se que o padrão de subida e descida se repete a cada 12 horas (maré semidiurna). Os dados coletados foram:

\begin{center}
\begin{tabular}{|p{0.143\linewidth}|p{0.143\linewidth}|p{0.143\linewidth}|p{0.143\linewidth}|p{0.143\linewidth}|p{0.143\linewidth}|}
\hline
\textbf{t (horas)} & \textbf{0} & \textbf{3} & \textbf{6} & \textbf{9} & \textbf{12} \\
\hline
h (metros) & 8 & 5 & 2 & 5 & 8 \\
\hline
\end{tabular}
\end{center}

Esse tipo de fenômeno periódico pode ser representado no plano cartesiano por uma função seno ou por uma função cosseno. Observando em que instante ocorrem os valores máximo e mínimo da maré, identifique qual das duas representações (seno ou cosseno) é a mais adequada para modelar h(t) e, em seguida, utilize o modelo escolhido para determinar a altura da água no instante t = 2 horas.
\begin{enumerate}[label=(\alph*),nosep,leftmargin=*]
\item 6,5 m
\item 7,6 m (aproximadamente)
\item 9,5 m
\item 3,5 m
\end{enumerate}

\textbf{Gabarito proposto:} 6,5 m

\textbf{Habilidade declarada:} EM13MAT306 --- Resolver e elaborar problemas em contextos que envolvem fenômenos periódicos reais, como ondas sonoras, ciclos menstruais, movimentos cíclicos, entre outros, e comparar suas representações com as funções seno e cosseno, no plano cartesiano, com ou sem apoio de aplicativos de álgebra e geometria.
\end{samepage}

\noindent\rule{\linewidth}{0.4pt}
\textbf{Tem erro matemático?}\quad
\mbox{$\square$~não} \quad \mbox{$\square$~sim, aqui:} \underline{\hspace{5cm}}
\quad \mbox{$\square$~não sei dizer}

\textbf{Que nível de dificuldade esta questão tem para os seus alunos?}\quad
\mbox{$\square$~fácil} \quad \mbox{$\square$~média} \quad
\mbox{$\square$~difícil} \quad \mbox{$\square$~fora do alcance da turma}

\textbf{Você usaria esta questão?}\quad
\mbox{$\square$~aceita} \quad \mbox{$\square$~aceita com ajuste} \quad
\mbox{$\square$~recusada}

\textbf{Por quê?} \emph{(obrigatório quando não for ``aceita'')}

\underline{\hspace{\linewidth}}

\underline{\hspace{\linewidth}}
\vspace{0.6em}

\begin{samepage}
\subsection*{Q042}
Uma equipe de engenharia rodoviária testa, em pista seca, a distância de frenagem d (em metros) de um automóvel a partir do instante em que o motorista aciona os freios, para diferentes velocidades v (em km/h). Os resultados de quatro testes estão na tabela abaixo:

\begin{center}
\begin{tabular}{|p{0.172\linewidth}|p{0.172\linewidth}|p{0.172\linewidth}|p{0.172\linewidth}|p{0.172\linewidth}|}
\hline
\textbf{v (km/h)} & \textbf{10} & \textbf{20} & \textbf{30} & \textbf{40} \\
\hline
d (m) & 2 & 8 & 18 & 32 \\
\hline
\end{tabular}
\end{center}

Sabendo que a distância de frenagem depende apenas da velocidade e que o padrão observado nos testes se mantém para qualquer velocidade, qual expressão algébrica representa corretamente d em função de v?
\begin{enumerate}[label=(\alph*),nosep,leftmargin=*]
\item $d(v) = \dfrac{v^2}{50}$
\item $d(v) = 0{,}6v - 4$
\item $d(v) = 0{,}2v$
\item $d(v) = \dfrac{v^2}{25}$
\end{enumerate}

\textbf{Gabarito proposto:} d(v) = v²/50

\textbf{Habilidade declarada:} EM13MAT502 --- Investigar relações entre números expressos em tabelas para representá-los no plano cartesiano, identificando padrões e criando conjecturas para generalizar e expressar algebricamente essa generalização, reconhecendo quando essa representação é de função polinomial de 2º grau do tipo y = ax².
\end{samepage}

\noindent\rule{\linewidth}{0.4pt}
\textbf{Tem erro matemático?}\quad
\mbox{$\square$~não} \quad \mbox{$\square$~sim, aqui:} \underline{\hspace{5cm}}
\quad \mbox{$\square$~não sei dizer}

\textbf{Que nível de dificuldade esta questão tem para os seus alunos?}\quad
\mbox{$\square$~fácil} \quad \mbox{$\square$~média} \quad
\mbox{$\square$~difícil} \quad \mbox{$\square$~fora do alcance da turma}

\textbf{Você usaria esta questão?}\quad
\mbox{$\square$~aceita} \quad \mbox{$\square$~aceita com ajuste} \quad
\mbox{$\square$~recusada}

\textbf{Por quê?} \emph{(obrigatório quando não for ``aceita'')}

\underline{\hspace{\linewidth}}

\underline{\hspace{\linewidth}}
\vspace{0.6em}

\begin{samepage}
\subsection*{Q025}
Considere as funções reais f(x) = 2x, g(x) = 2x + 3 e h(x) = -2x, definidas para todo x real.

a) Determine as coordenadas de dois pontos do gráfico de cada função e utilize-os para esboçar, num mesmo plano cartesiano, as retas correspondentes a f, g e h.

b) Observando os três gráficos esboçados, indique quais retas passam pela origem do plano cartesiano e quais não passam.

c) A partir do que você observou, enuncie uma regra geral: dada uma lei y = ax + b (com a $\neq$ 0), como decidir, apenas olhando os coeficientes a e b (sem construir o gráfico), se a reta correspondente passará ou não pela origem? Qual é o nome usual dado ao tipo de função cujo gráfico passa pela origem, e qual o nome dado ao caso em que isso não ocorre?

d) Compare os coeficientes angulares de f, g e h. Existe alguma relação geométrica de paralelismo entre duas dessas retas? Justifique sua resposta e descreva, em palavras, como o gráfico de g pode ser obtido a partir do gráfico de f.

\textbf{Gabarito proposto:} b) f e h passam pela origem; g não passa. c) A reta y=ax+b passa pela origem se, e somente se, b=0; nesse caso a função é linear/proporcional, e quando b$\neq$0 é afim não proporcional. d) f e g são paralelas (mesmo a=2); h não é paralela a nenhuma delas; o gráfico de g é o gráfico de f transladado 3 unidades para cima.

\textbf{Habilidade declarada:} EM13MAT401 --- Converter representações algébricas de funções polinomiais de 1º grau para representações geométricas no plano cartesiano, distinguindo os casos nos quais o comportamento é proporcional, recorrendo ou não a softwares ou aplicativos de álgebra e geometria dinâmica.
\end{samepage}

\noindent\rule{\linewidth}{0.4pt}
\textbf{Tem erro matemático?}\quad
\mbox{$\square$~não} \quad \mbox{$\square$~sim, aqui:} \underline{\hspace{5cm}}
\quad \mbox{$\square$~não sei dizer}

\textbf{Que nível de dificuldade esta questão tem para os seus alunos?}\quad
\mbox{$\square$~fácil} \quad \mbox{$\square$~média} \quad
\mbox{$\square$~difícil} \quad \mbox{$\square$~fora do alcance da turma}

\textbf{Você usaria esta questão?}\quad
\mbox{$\square$~aceita} \quad \mbox{$\square$~aceita com ajuste} \quad
\mbox{$\square$~recusada}

\textbf{Por quê?} \emph{(obrigatório quando não for ``aceita'')}

\underline{\hspace{\linewidth}}

\underline{\hspace{\linewidth}}
\vspace{0.6em}

\begin{samepage}
\subsection*{Q034}
Duas empresas de aplicativo de transporte, A e B, calculam o valor da corrida em função da distância percorrida $x$ (em km). A empresa A cobra R\$\,2,50 por quilômetro rodado, sem nenhuma taxa adicional. A empresa B cobra uma taxa fixa de embarque de R\$\,3,00 mais R\$\,2,00 por quilômetro rodado. Sejam $y_A(x)$ e $y_B(x)$, em reais, os custos totais das corridas das empresas A e B, respectivamente, em função de $x$. Se essas duas funções forem representadas por retas em um plano cartesiano (custo $y$ no eixo vertical e distância $x$ no eixo horizontal), assinale a alternativa que descreve corretamente essas retas.
\begin{enumerate}[label=(\alph*),nosep,leftmargin=*]
\item A reta que representa a empresa A passa pela origem (0,0) e tem coeficiente angular 2,5, sendo uma função proporcional; a reta da empresa B corta o eixo y no ponto (0,3) e tem coeficiente angular 2, não sendo proporcional.
\item A reta que representa a empresa B passa pela origem (0,0), pois seu coeficiente angular é 2; a reta da empresa A corta o eixo y no ponto (0,3), pois seu coeficiente angular é 2,5.
\item Ambas as retas passam pela origem (0,0), já que representam custo por quilômetro rodado; a reta de A tem coeficiente angular 2,5 e a de B tem coeficiente angular 2.
\item A reta da empresa A corta o eixo y no ponto (0, 2,5); a reta da empresa B corta o eixo y no ponto (0,2) e tem coeficiente angular 3.
\end{enumerate}

\textbf{Gabarito proposto:} A reta que representa a empresa A passa pela origem $(0,0)$ e tem coeficiente angular 2,5 (função proporcional); a reta da empresa B corta o eixo y no ponto $(0,3)$ e tem coeficiente angular 2 (função afim, não proporcional).

\textbf{Habilidade declarada:} EM13MAT401 --- Converter representações algébricas de funções polinomiais de 1º grau para representações geométricas no plano cartesiano, distinguindo os casos nos quais o comportamento é proporcional, recorrendo ou não a softwares ou aplicativos de álgebra e geometria dinâmica.
\end{samepage}

\noindent\rule{\linewidth}{0.4pt}
\textbf{Tem erro matemático?}\quad
\mbox{$\square$~não} \quad \mbox{$\square$~sim, aqui:} \underline{\hspace{5cm}}
\quad \mbox{$\square$~não sei dizer}

\textbf{Que nível de dificuldade esta questão tem para os seus alunos?}\quad
\mbox{$\square$~fácil} \quad \mbox{$\square$~média} \quad
\mbox{$\square$~difícil} \quad \mbox{$\square$~fora do alcance da turma}

\textbf{Você usaria esta questão?}\quad
\mbox{$\square$~aceita} \quad \mbox{$\square$~aceita com ajuste} \quad
\mbox{$\square$~recusada}

\textbf{Por quê?} \emph{(obrigatório quando não for ``aceita'')}

\underline{\hspace{\linewidth}}

\underline{\hspace{\linewidth}}
\vspace{0.6em}

\begin{samepage}
\subsection*{Q004}
Uma companhia de saneamento calcula o valor mensal da conta de água de acordo com o consumo $x$ (em m³) do cliente, segundo a tabela abaixo:

\begin{center}
\begin{tabular}{|p{0.430\linewidth}|p{0.430\linewidth}|}
\hline
\textbf{Consumo mensal} & \textbf{Regra de cobrança} \\
\hline
$0 \le x \le 10$ & taxa mínima fixa de R\$\,20,00 \\
\hline
$10 < x \le 20$ & R\$\,20,00 mais R\$\,3,00 por cada m³ que exceder 10 \\
\hline
$x > 20$ & R\$\,50,00 mais R\$\,5,00 por cada m³ que exceder 20 \\
\hline
\end{tabular}
\end{center}

Seja $f(x)$ o valor, em reais, da conta de água correspondente a um consumo de $x$ m³.

a) Escreva a lei de formação algébrica de $f(x)$, especificando a expressão válida em cada uma das três faixas de consumo.

b) Qual é o domínio de validade de $f$, isto é, o conjunto de valores de $x$ para os quais a regra de cobrança faz sentido?

c) Uma família definiu que pode gastar, no máximo, R\$\,41,00 por mês com água, e sabe que seu consumo está entre 10 m³ e 20 m³. Qual é o maior consumo mensal, em m³, que essa família pode ter sem ultrapassar esse orçamento?

d) Esboce o gráfico de $f$ para $0 \le x \le 30$ e classifique o comportamento de $f$ (crescente, decrescente ou constante) em cada um dos três trechos.

e) Determine a imagem de $f$ restrita ao intervalo $0 \le x \le 30$.

\textbf{Gabarito proposto:} a) $f(x)=20$ se $0\le x\le 10$; $f(x)=3x-10$ se $10<x\le 20$; $f(x)=5x-50$ se $x>20$. b) $D(f)=[0,+\infty)$. c) 17 m³. d) constante em $[0,10]$; crescente em $(10,20]$ e em $(20,30]$ (com inclinação maior nesse último trecho). e) Imagem $=[20,100]$.

\textbf{Habilidade declarada:} EM13MAT405 --- Reconhecer funções definidas por uma ou mais sentenças (como a tabela do Imposto de Renda, contas de luz, água, gás etc.), em suas representações algébrica e gráfica, convertendo essas representações de uma para outra e identificando domínios de validade, imagem, crescimento e decrescimento.
\end{samepage}

\noindent\rule{\linewidth}{0.4pt}
\textbf{Tem erro matemático?}\quad
\mbox{$\square$~não} \quad \mbox{$\square$~sim, aqui:} \underline{\hspace{5cm}}
\quad \mbox{$\square$~não sei dizer}

\textbf{Que nível de dificuldade esta questão tem para os seus alunos?}\quad
\mbox{$\square$~fácil} \quad \mbox{$\square$~média} \quad
\mbox{$\square$~difícil} \quad \mbox{$\square$~fora do alcance da turma}

\textbf{Você usaria esta questão?}\quad
\mbox{$\square$~aceita} \quad \mbox{$\square$~aceita com ajuste} \quad
\mbox{$\square$~recusada}

\textbf{Por quê?} \emph{(obrigatório quando não for ``aceita'')}

\underline{\hspace{\linewidth}}

\underline{\hspace{\linewidth}}
\vspace{0.6em}

\begin{samepage}
\subsection*{Q087}
Considere a progressão aritmética $(a_n)$ com primeiro termo $a_1 = 7$ e razão $r = 4$: $(7, 11, 15, 19, 23, \dots)$. Essa PA pode ser vista como a restrição da função afim $f(x) = 4x + 3$ ao conjunto dos números naturais não nulos, de modo que $a_n = f(n)$ para todo $n \geq 1$ (o gráfico da PA é o conjunto de pontos discretos $(n, a_n)$ sobre a reta que representa $f$).

Qual é o menor termo dessa PA que é maior do que 100?
\begin{enumerate}[label=(\alph*),nosep,leftmargin=*]
\item 103
\item 99
\item 107
\item 100
\end{enumerate}

\textbf{Gabarito proposto:} 103

\textbf{Habilidade declarada:} EM13MAT507 --- Identificar e associar sequências numéricas (PA) a funções afins de domínios discretos para análise de propriedades, incluindo dedução de algumas fórmulas e resolução de problemas.
\end{samepage}

\noindent\rule{\linewidth}{0.4pt}
\textbf{Tem erro matemático?}\quad
\mbox{$\square$~não} \quad \mbox{$\square$~sim, aqui:} \underline{\hspace{5cm}}
\quad \mbox{$\square$~não sei dizer}

\textbf{Que nível de dificuldade esta questão tem para os seus alunos?}\quad
\mbox{$\square$~fácil} \quad \mbox{$\square$~média} \quad
\mbox{$\square$~difícil} \quad \mbox{$\square$~fora do alcance da turma}

\textbf{Você usaria esta questão?}\quad
\mbox{$\square$~aceita} \quad \mbox{$\square$~aceita com ajuste} \quad
\mbox{$\square$~recusada}

\textbf{Por quê?} \emph{(obrigatório quando não for ``aceita'')}

\underline{\hspace{\linewidth}}

\underline{\hspace{\linewidth}}
\vspace{0.6em}

\begin{samepage}
\subsection*{Q077}
Uma fintech oferece uma aplicação financeira em que o capital investido é corrigido mensalmente por juros compostos, a uma taxa fixa, sem aportes ou resgates. Marina investiu R\$\,6.000,00 nessa aplicação, com taxa de 4\% ao mês.

a) Elabore a expressão algébrica da função $M(t)$, que fornece o montante (em reais) acumulado por Marina após $t$ meses de aplicação.

b) Utilizando a função elaborada, calcule o montante ao final do 10º mês, arredondando o resultado para duas casas decimais.

c) Determine o menor número inteiro de meses necessário para que o montante ultrapasse o dobro do capital investido. Em seguida, explique por que esse tempo é bem menor do que o que seria necessário caso o capital rendesse mensalmente um valor fixo de R\$\,240,00 (4\% do capital inicial, sem incidência de juros sobre juros) — evidenciando, com essa comparação, o caráter exponencial (e não linear) do crescimento sob juros compostos.

\textbf{Gabarito proposto:} a) $M(t) = 6000\cdot(1,04)^t$; b) $M(10) \approx R\$\,8.881{,}47$; c) menor número inteiro de meses: $t = 18$ (contra 25 meses no regime linear equivalente, evidenciando o crescimento exponencial).

\textbf{Habilidade declarada:} EM13MAT303 --- Resolver e elaborar problemas envolvendo porcentagens em diversos contextos e sobre juros compostos, destacando o crescimento exponencial.
\end{samepage}

\noindent\rule{\linewidth}{0.4pt}
\textbf{Tem erro matemático?}\quad
\mbox{$\square$~não} \quad \mbox{$\square$~sim, aqui:} \underline{\hspace{5cm}}
\quad \mbox{$\square$~não sei dizer}

\textbf{Que nível de dificuldade esta questão tem para os seus alunos?}\quad
\mbox{$\square$~fácil} \quad \mbox{$\square$~média} \quad
\mbox{$\square$~difícil} \quad \mbox{$\square$~fora do alcance da turma}

\textbf{Você usaria esta questão?}\quad
\mbox{$\square$~aceita} \quad \mbox{$\square$~aceita com ajuste} \quad
\mbox{$\square$~recusada}

\textbf{Por quê?} \emph{(obrigatório quando não for ``aceita'')}

\underline{\hspace{\linewidth}}

\underline{\hspace{\linewidth}}
\vspace{0.6em}

\begin{samepage}
\subsection*{Q024}
A tabela abaixo relaciona valores de $x$ e os correspondentes valores de $y$, obtidos a partir de uma mesma regra de formação (os valores de $x$ não variam em intervalos iguais):

\begin{center}
\begin{tabular}{|p{0.172\linewidth}|p{0.172\linewidth}|p{0.172\linewidth}|p{0.172\linewidth}|p{0.172\linewidth}|}
\hline
\textbf{$x$} & \textbf{$-3$} & \textbf{$1$} & \textbf{$4$} & \textbf{$8$} \\
\hline
$y$ & $14$ & $2$ & $-7$ & $-19$ \\
\hline
\end{tabular}
\end{center}

Um estudante, analisando os pares ordenados $(x,y)$ da tabela, conjectura que essa relação pode ser representada por uma função polinomial do 1º grau. Assinale a alternativa que apresenta a lei de formação $f(x)$ que generaliza corretamente o padrão observado na tabela.
\begin{enumerate}[label=(\alph*),nosep,leftmargin=*]
\item $f(x) = -3x + 5$
\item $f(x) = -12x + 14$
\item $f(x) = 3x - 1$
\item $f(x) = -11x + 13$
\end{enumerate}

\textbf{Gabarito proposto:} f(x) = -3x + 5

\textbf{Habilidade declarada:} EM13MAT501 --- Investigar relações entre números expressos em tabelas para representá-los no plano cartesiano, identificando padrões e criando conjecturas para generalizar e expressar algebricamente essa generalização, reconhecendo quando essa representação é de função polinomial de 1º grau.
\end{samepage}

\noindent\rule{\linewidth}{0.4pt}
\textbf{Tem erro matemático?}\quad
\mbox{$\square$~não} \quad \mbox{$\square$~sim, aqui:} \underline{\hspace{5cm}}
\quad \mbox{$\square$~não sei dizer}

\textbf{Que nível de dificuldade esta questão tem para os seus alunos?}\quad
\mbox{$\square$~fácil} \quad \mbox{$\square$~média} \quad
\mbox{$\square$~difícil} \quad \mbox{$\square$~fora do alcance da turma}

\textbf{Você usaria esta questão?}\quad
\mbox{$\square$~aceita} \quad \mbox{$\square$~aceita com ajuste} \quad
\mbox{$\square$~recusada}

\textbf{Por quê?} \emph{(obrigatório quando não for ``aceita'')}

\underline{\hspace{\linewidth}}

\underline{\hspace{\linewidth}}
\vspace{0.6em}

\begin{samepage}
\subsection*{Q039}
A tabela abaixo relaciona o comprimento $L$ do lado de um certo tipo de figura com a grandeza $A$ associada a ela, para quatro valores de $L$:

\begin{center}
\begin{tabular}{|p{0.430\linewidth}|p{0.430\linewidth}|}
\hline
\textbf{$L$} & \textbf{$A$} \\
\hline
1 & 3 \\
\hline
2 & 12 \\
\hline
3 & 27 \\
\hline
4 & 48 \\
\hline
\end{tabular}
\end{center}

Investigando os pares de valores dessa tabela, qual expressão algébrica, escrita em função de $L$, generaliza corretamente a relação entre $A$ e $L$ para todos os pares apresentados (e não apenas para alguns deles)?
\begin{enumerate}[label=(\alph*),nosep,leftmargin=*]
\item $A = 3L^2$
\item $A = 3L$
\item $A = L^2 + 2$
\item $A = 9L - 6$
\end{enumerate}

\textbf{Gabarito proposto:} A = 3L²

\textbf{Habilidade declarada:} EM13MAT502 --- Investigar relações entre números expressos em tabelas para representá-los no plano cartesiano, identificando padrões e criando conjecturas para generalizar e expressar algebricamente essa generalização, reconhecendo quando essa representação é de função polinomial de 2º grau do tipo y = ax².
\end{samepage}

\noindent\rule{\linewidth}{0.4pt}
\textbf{Tem erro matemático?}\quad
\mbox{$\square$~não} \quad \mbox{$\square$~sim, aqui:} \underline{\hspace{5cm}}
\quad \mbox{$\square$~não sei dizer}

\textbf{Que nível de dificuldade esta questão tem para os seus alunos?}\quad
\mbox{$\square$~fácil} \quad \mbox{$\square$~média} \quad
\mbox{$\square$~difícil} \quad \mbox{$\square$~fora do alcance da turma}

\textbf{Você usaria esta questão?}\quad
\mbox{$\square$~aceita} \quad \mbox{$\square$~aceita com ajuste} \quad
\mbox{$\square$~recusada}

\textbf{Por quê?} \emph{(obrigatório quando não for ``aceita'')}

\underline{\hspace{\linewidth}}

\underline{\hspace{\linewidth}}
\vspace{0.6em}

\begin{samepage}
\subsection*{Q065}
Um oceanógrafo registra a altura da maré (em metros) em um pequeno porto ao longo de um dia. Ele observa que a maré varia de forma periódica e simétrica em torno de um nível médio, atingindo seu valor máximo de 1,8 m exatamente à meia-noite ($t=0$, com $t$ em horas) e seu valor mínimo de 0,2 m exatamente às 6h da manhã. Esse padrão se repete a cada 12 horas.

a) Modele a altura da maré, em metros, por uma função do tipo $h(t) = A\cos(Bt) + D$, com $t$ em horas contado a partir da meia-noite. Determine explicitamente os valores de $A$, $B$ e $D$, justificando cada um a partir dos dados do fenômeno (amplitude, período e deslocamento vertical do gráfico em relação ao eixo $t$).

b) Determine todos os horários dentro das primeiras 12 horas do dia (isto é, para $0 \le t < 12$) em que a maré atinge exatamente 1,4 m de altura. Para cada horário encontrado, indique se a maré está subindo ou descendo naquele instante, justificando sua resposta apenas a partir do comportamento gráfico da função cosseno (sem usar derivadas).

\textbf{Gabarito proposto:} h(t) = 0,8·cos($\pi$t/6) + 1,0 (A=0,8 m; B=$\pi$/6 rad/h; D=1,0 m). A maré atinge 1,4 m às 2h (descendo) e às 10h (subindo).

\textbf{Habilidade declarada:} EM13MAT306 --- Resolver e elaborar problemas em contextos que envolvem fenômenos periódicos reais, como ondas sonoras, ciclos menstruais, movimentos cíclicos, entre outros, e comparar suas representações com as funções seno e cosseno, no plano cartesiano, com ou sem apoio de aplicativos de álgebra e geometria.
\end{samepage}

\noindent\rule{\linewidth}{0.4pt}
\textbf{Tem erro matemático?}\quad
\mbox{$\square$~não} \quad \mbox{$\square$~sim, aqui:} \underline{\hspace{5cm}}
\quad \mbox{$\square$~não sei dizer}

\textbf{Que nível de dificuldade esta questão tem para os seus alunos?}\quad
\mbox{$\square$~fácil} \quad \mbox{$\square$~média} \quad
\mbox{$\square$~difícil} \quad \mbox{$\square$~fora do alcance da turma}

\textbf{Você usaria esta questão?}\quad
\mbox{$\square$~aceita} \quad \mbox{$\square$~aceita com ajuste} \quad
\mbox{$\square$~recusada}

\textbf{Por quê?} \emph{(obrigatório quando não for ``aceita'')}

\underline{\hspace{\linewidth}}

\underline{\hspace{\linewidth}}
\vspace{0.6em}

\begin{samepage}
\subsection*{Q015}
Uma colônia de bactérias em um experimento de laboratório tem seu número de indivíduos, a partir do início da observação, dado por $N(t) = 100 \cdot 2^{t}$, em que $t$ é o tempo em horas ($t \geq 0$) e $N(t)$ é o número de bactérias. Um pesquisador quer saber, ao contrário, quanto tempo é necessário para que a colônia atinja um determinado número $N$ de bactérias, e por isso utiliza a função inversa de $N(t)$, que ele chama de $t(N)$.

a) Determine o domínio e a imagem de $N(t)$, considerando o contexto do experimento.

b) Encontre a expressão de $t(N)$ (a função inversa de $N(t)$) e determine seu domínio e sua imagem.

c) Compare o crescimento das duas funções: as duas são crescentes, mas uma cresce 'cada vez mais rápido' e a outra cresce 'cada vez mais devagar' à medida que a variável aumenta. Identifique qual é qual e justifique observando como varia o incremento de cada função.

d) Explique por que o domínio de $N(t)$ coincide com a imagem de $t(N)$, e a imagem de $N(t)$ coincide com o domínio de $t(N)$, relacionando esse fato com a posição dos gráficos das duas funções no plano cartesiano.

\textbf{Gabarito proposto:} a) $D_N=[0,+\infty)$, $Im_N=[100,+\infty)$. b) $t(N)=\log_2(N/100)$, com $D_t=[100,+\infty)$, $Im_t=[0,+\infty)$. c) $N(t)$ cresce cada vez mais rápido (convexa); $t(N)$ cresce cada vez mais devagar (côncava). d) Por serem funções inversas, seus gráficos são simétricos em relação à reta $y=x$, o que troca domínio e imagem entre as duas funções.

\textbf{Habilidade declarada:} EM13MAT403 --- Comparar e analisar as representações, em plano cartesiano, das funções exponencial e logarítmica para identificar as características fundamentais (domínio, imagem, crescimento) de cada uma, com ou sem apoio de tecnologias digitais, estabelecendo relações entre elas.
\end{samepage}

\noindent\rule{\linewidth}{0.4pt}
\textbf{Tem erro matemático?}\quad
\mbox{$\square$~não} \quad \mbox{$\square$~sim, aqui:} \underline{\hspace{5cm}}
\quad \mbox{$\square$~não sei dizer}

\textbf{Que nível de dificuldade esta questão tem para os seus alunos?}\quad
\mbox{$\square$~fácil} \quad \mbox{$\square$~média} \quad
\mbox{$\square$~difícil} \quad \mbox{$\square$~fora do alcance da turma}

\textbf{Você usaria esta questão?}\quad
\mbox{$\square$~aceita} \quad \mbox{$\square$~aceita com ajuste} \quad
\mbox{$\square$~recusada}

\textbf{Por quê?} \emph{(obrigatório quando não for ``aceita'')}

\underline{\hspace{\linewidth}}

\underline{\hspace{\linewidth}}
\vspace{0.6em}

\begin{samepage}
\subsection*{Q085}
Um capital $C_0$ é aplicado a juros compostos, com taxa fixa $i$ por período, de modo que o montante após $n$ períodos é dado por $M(n) = C_0\cdot(1+i)^n$. Sabe-se que, ao final de 3 períodos, o montante corresponde a 8 vezes o capital inicial, isto é, $M(3) = 8\,C_0$. Mantendo-se a mesma taxa $i$ por todo o tempo, qual é o montante $M(6)$, em função de $C_0$?
\begin{enumerate}[label=(\alph*),nosep,leftmargin=*]
\item $16\,C_0$
\item $32\,C_0$
\item $64\,C_0$
\item $512\,C_0$
\end{enumerate}

\textbf{Gabarito proposto:} 64 C0

\textbf{Habilidade declarada:} EM13MAT303 --- Resolver e elaborar problemas envolvendo porcentagens em diversos contextos e sobre juros compostos, destacando o crescimento exponencial.
\end{samepage}

\noindent\rule{\linewidth}{0.4pt}
\textbf{Tem erro matemático?}\quad
\mbox{$\square$~não} \quad \mbox{$\square$~sim, aqui:} \underline{\hspace{5cm}}
\quad \mbox{$\square$~não sei dizer}

\textbf{Que nível de dificuldade esta questão tem para os seus alunos?}\quad
\mbox{$\square$~fácil} \quad \mbox{$\square$~média} \quad
\mbox{$\square$~difícil} \quad \mbox{$\square$~fora do alcance da turma}

\textbf{Você usaria esta questão?}\quad
\mbox{$\square$~aceita} \quad \mbox{$\square$~aceita com ajuste} \quad
\mbox{$\square$~recusada}

\textbf{Por quê?} \emph{(obrigatório quando não for ``aceita'')}

\underline{\hspace{\linewidth}}

\underline{\hspace{\linewidth}}
\vspace{0.6em}

\begin{samepage}
\subsection*{Q086}
Uma roda-gigante tem 40 m de diâmetro e seu centro está a 22 m de altura em relação ao solo. Ela gira com velocidade constante, completando uma volta a cada 12 minutos. Um passageiro embarca no ponto mais baixo da roda (rente ao chão de segurança, a 2 m de altura) no instante $t = 0$ minuto. A altura $h(t)$, em metros, desse passageiro em relação ao solo, $t$ minutos após o embarque, pode ser descrita por uma função do tipo cosseno. Considerando o gráfico de $h(t)$ ao longo de todo o primeiro giro completo ($0 \le t \le 12$ min), assinale a alternativa que apresenta corretamente a lei de $h(t)$ e o intervalo de tempo, nesse primeiro giro, em que a altura do passageiro é maior que 32 m.
\begin{enumerate}[label=(\alph*),nosep,leftmargin=*]
\item $h(t) = 22 - 20\cos\left(\dfrac{\pi t}{6}\right)$, e a altura é maior que 32 m para $t \in (4, 8)$ minutos.
\item $h(t) = 22 - 40\cos\left(\dfrac{\pi t}{6}\right)$, e a altura é maior que 32 m para $t \in (4, 8)$ minutos.
\item $h(t) = 22 - 20\cos\left(\dfrac{\pi t}{3}\right)$, e a altura é maior que 32 m para $t \in (2, 4)$ minutos.
\item $h(t) = 22 + 20\cos\left(\dfrac{\pi t}{6}\right)$, e a altura é maior que 32 m para $t \in (0, 2) \cup (10, 12)$ minutos.
\end{enumerate}

\textbf{Gabarito proposto:} h(t) = 22 - 20cos($\pi$t/6); altura > 32 m para t entre 4 e 8 minutos (alternativa A)

\textbf{Habilidade declarada:} EM13MAT306 --- Resolver e elaborar problemas em contextos que envolvem fenômenos periódicos reais, como ondas sonoras, ciclos menstruais, movimentos cíclicos, entre outros, e comparar suas representações com as funções seno e cosseno, no plano cartesiano, com ou sem apoio de aplicativos de álgebra e geometria.
\end{samepage}

\noindent\rule{\linewidth}{0.4pt}
\textbf{Tem erro matemático?}\quad
\mbox{$\square$~não} \quad \mbox{$\square$~sim, aqui:} \underline{\hspace{5cm}}
\quad \mbox{$\square$~não sei dizer}

\textbf{Que nível de dificuldade esta questão tem para os seus alunos?}\quad
\mbox{$\square$~fácil} \quad \mbox{$\square$~média} \quad
\mbox{$\square$~difícil} \quad \mbox{$\square$~fora do alcance da turma}

\textbf{Você usaria esta questão?}\quad
\mbox{$\square$~aceita} \quad \mbox{$\square$~aceita com ajuste} \quad
\mbox{$\square$~recusada}

\textbf{Por quê?} \emph{(obrigatório quando não for ``aceita'')}

\underline{\hspace{\linewidth}}

\underline{\hspace{\linewidth}}
\vspace{0.6em}

\newpage
\subsection*{Para terminar}

\textbf{1. Você usaria um sistema assim no seu planejamento? Por quê?}

\underline{\hspace{\linewidth}}

\underline{\hspace{\linewidth}}

\bigskip
\textbf{2. O que faltou nas questões que você viu?}

\underline{\hspace{\linewidth}}

\underline{\hspace{\linewidth}}

\bigskip
\textbf{3. O que atrapalhou --- na questão, no enunciado, no formato deste
material?}

\underline{\hspace{\linewidth}}

\underline{\hspace{\linewidth}}

\bigskip
Obrigado. Se quiser saber quais itens tinham erro e qual era, é só pedir.
\end{document}